\documentclass[12pt,letterpaper]{article}

% --- Paquetes Requeridos ---
\usepackage[utf8]{inputenc}
\usepackage[spanish,es-tabla]{babel}
\usepackage[margin=2.5cm]{geometry}
\usepackage{graphicx}
\usepackage{float}
\usepackage{titlesec}
\usepackage{setspace}
\usepackage{hyperref}
\usepackage{booktabs}
\usepackage{caption}
\usepackage{enumitem}
\usepackage{listings}
\usepackage{xcolor}

% --- Configuración de Párrafo ---
\onehalfspacing
\setlength{\parindent}{0pt}
\setlength{\parskip}{1em}

% --- Configuración de listings ---
\lstset{
    basicstyle=\ttfamily\footnotesize,
    breaklines=true,
    frame=single,
    captionpos=b,
    showstringspaces=false
}

% --- Inicio del Documento ---
\begin{document}

% --- PORTADA ---
\begin{titlepage}
    \centering

    {\Large \textbf{UNIVERSIDAD SANTO TOMÁS}} \\[1cm]

    {\Large \textbf{Sistema colaborativo de robots móviles para asistencia de orientación en entornos interiores con múltiples pisos}} \\[2cm]

    {\large \textbf{Realizado por:}} \\
    {\large Santiago Hernández Ávila} \\
    {\large Jonny Alejandro Mejía León} \\[1cm]

    \begin{figure}[H]
        \centering
        \includegraphics[width=0.6\textwidth]{Logos Uni.png}
        \label{Logo Universidad Santo Tomas y facultad de ingenieria electronica}
    \end{figure}

    \vfill

    {\large \textbf{Cronograma de actividades --- Semanas 17 a 32}} \\
    {\large Proyecto de Grado 2 --- Fase 5: Integración del sistema y realización de pruebas} \\[1.5cm]

    \textbf{Dirigido por:} \\
    Ing. Armando Mateus Rojas, Msc. \\
    Ing. Nestor Ivan Ospina, Msc. \\
    Ing. Oscar Mauricio Gélvez Lizarazo, Msc. \\[1.0cm]

    \vfill

    {\large GED (Grupo de Estudio y Desarrollo en Robótica)} \\
    {\large Facultad de Ingeniería Electrónica} \\
    {\large División de Ingenierías} \\
    {\large Bogotá, agosto de 2026}
\end{titlepage}

% --- CONTENIDO ---

\section{Presentación}

El presente documento desarrolla el cronograma de actividades del proyecto para el segundo periodo académico, correspondiente a las semanas S17 a S32. Las actividades aquí programadas son las contempladas en el Capítulo 8 del anteproyecto \cite{anteproyecto}; el documento las distribuye semana a semana y les asigna un criterio de cierre verificable.

El proyecto se desarrolla en dos periodos académicos separados por el receso de mitad de año, según se indica en la Tabla \ref{tab:periodos}. La programación de las semanas S1 a S16 se ejecutó durante el primer periodo; este documento cubre el segundo.

\begin{table}[H]
    \centering
    \caption{Distribución del proyecto en los dos periodos académicos.}
    \label{tab:periodos}
    \begin{tabular}{@{}lll@{}}
        \toprule
        \textbf{Periodo} & \textbf{Semanas} & \textbf{Calendario} \\ \midrule
        Proyecto de Grado 1 & S1--S16  & febrero -- mayo de 2026 \\
        Proyecto de Grado 2 & S17--S32 & agosto -- noviembre de 2026 \\ \bottomrule
    \end{tabular}
\end{table}

Sobre la marcha del desarrollo se precisaron las ventanas temporales de algunas actividades, conforme se detalla en la Sección \ref{sec:ajustes}. Se trata de ajustes propios de la ejecución del proyecto: todas las actividades del anteproyecto se conservan, con su mismo alcance y su misma secuencia de dependencias.

\section{Calendario de referencia}

\begin{table}[H]
    \centering
    \caption{Calendario de las semanas S17 a S32.}
    \label{tab:calendario}
    \begin{tabular}{@{}llll@{}}
        \toprule
        \textbf{Semana} & \textbf{Fechas} & \textbf{Semana} & \textbf{Fechas} \\ \midrule
        S17 & 3 -- 9 ago          & S25 & 28 sep -- 4 oct \\
        S18 & 10 -- 16 ago        & S26 & 5 -- 11 oct \\
        S19 & 17 -- 23 ago        & S27 & 12 -- 18 oct \\
        S20 & 24 -- 30 ago        & \textbf{S28} & \textbf{19 -- 25 oct} \\
        S21 & 31 ago -- 6 sep     & \textbf{S29} & \textbf{26 oct -- 1 nov} \\
        S22 & 7 -- 13 sep         & S30 & 2 -- 8 nov \\
        S23 & 14 -- 20 sep        & S31 & 9 -- 15 nov \\
        S24 & 21 -- 27 sep        & S32 & 16 -- 22 nov \\ \bottomrule
    \end{tabular}
\end{table}

La sustentación se prevé a finales de octubre, esto es, en las semanas \textbf{S28--S29}. El documento final se entrega con posterioridad a la sustentación, conforme a la norma del espacio académico, de modo que incorpore las observaciones del jurado.

\subsection{Días hábiles disponibles}

La programación se expresa en semanas, pero el trabajo se ejecuta en días hábiles. Dentro del periodo S17--S32 hay cinco festivos, de modo que el presupuesto real es de \textbf{75 días hábiles} y no de 80. Se relacionan a continuación porque tres de ellos caen en semanas con criterio de cierre comprometido, y uno en la semana de entrega final.

\begin{table}[H]
    \centering
    \caption{Festivos comprendidos entre las semanas S17 y S32.}
    \label{tab:festivos}
    \begin{tabular}{@{}llc@{}}
        \toprule
        \textbf{Fecha} & \textbf{Festivo} & \textbf{Semana} \\ \midrule
        viernes 7 de agosto    & Batalla de Boyacá                & S17 \\
        lunes 17 de agosto     & Asunción de la Virgen            & S19 \\
        lunes 12 de octubre    & Día de la Raza                   & S27 \\
        lunes 2 de noviembre   & Todos los Santos                 & S30 \\
        lunes 16 de noviembre  & Independencia de Cartagena       & S32 \\ \bottomrule
    \end{tabular}
\end{table}

La consecuencia se refleja en la programación: S17 dispone de cuatro días hábiles, S18 de cinco y S19 de cuatro. Las actividades se distribuyeron sobre esa disponibilidad y no sobre semanas completas.

\section{Programación de actividades}

\subsection{Diagrama de programación}

\begin{table}[H]
    \centering
    \footnotesize
    \setlength{\tabcolsep}{2.4pt}
    \caption{Diagrama de programación de actividades para las semanas S17 a S32.}
    \label{tab:gantt}
    \begin{tabular}{@{}p{5.1cm}cccccccccccccccc@{}}
        \toprule
        \textbf{Actividad} & 17 & 18 & 19 & 20 & 21 & 22 & 23 & 24 & 25 & 26 & 27 & 28 & 29 & 30 & 31 & 32 \\ \midrule
        Ajuste progresivo de los elementos desarrollados & $\bullet$ & $\bullet$ & & & & & & & & & & & & & & \\
        Integración de los módulos del sistema & & $\bullet$ & $\bullet$ & $\bullet$ & $\bullet$ & $\bullet$ & & & & & & & & & & \\
        Verificación del funcionamiento general & & & & & & & $\bullet$ & & & & & & & & & \\
        Corrección de fallas o inconsistencias & & & & & & & $\bullet$ & & & & & & & & & \\
        Pruebas en condiciones controladas & & & & & & & $\bullet$ & $\bullet$ & $\bullet$ & & & & & & & \\
        Ejecución de pruebas de evaluación & & & & & & & & $\bullet$ & $\bullet$ & & & & & & & \\
        Recolección de información & & & & & & & & $\bullet$ & $\bullet$ & & & & & & & \\
        Análisis de resultados & & & & & & & & & & $\bullet$ & & & & & & \\
        Elaboración de conclusiones & & & & & & & & & & & $\bullet$ & & & & & \\
        Preparación de la sustentación final & & & & & & & & & & $\bullet$ & $\bullet$ & $\bullet$ & & & & \\
        Sustentación & & & & & & & & & & & & $\bullet$ & $\bullet$ & & & \\
        Elaboración del documento final & & & & & & & & & & & & & $\bullet$ & $\bullet$ & $\bullet$ & $\bullet$ \\
        Registro de avances del proyecto & $\bullet$ & $\bullet$ & $\bullet$ & $\bullet$ & $\bullet$ & $\bullet$ & $\bullet$ & $\bullet$ & $\bullet$ & $\bullet$ & $\bullet$ & $\bullet$ & $\bullet$ & $\bullet$ & $\bullet$ & $\bullet$ \\ \bottomrule
    \end{tabular}
\end{table}

\subsection{Ajustes de programación realizados durante la ejecución}
\label{sec:ajustes}

Las ventanas temporales de la Tabla \ref{tab:gantt} incorporan las precisiones que el desarrollo del proyecto hizo necesarias durante el primer periodo. Se enumeran a continuación:

\begin{itemize}
    \item La actividad \textit{Ajuste progresivo de los elementos desarrollados} se extiende hasta S18. La puesta a punto de la navegación autónoma de un agente exigió actualizar la pila de navegación a la distribución de ROS 2 empleada por el proyecto y adecuarla a la cinemática Ackermann del vehículo, tareas contempladas dentro de esta misma actividad.
    \item Las actividades que dependen de la anterior --- \textit{Integración de los módulos}, \textit{Verificación del funcionamiento general}, \textit{Pruebas en condiciones controladas}, \textit{Ejecución de pruebas de evaluación} y \textit{Análisis de resultados} --- se desplazan de forma consistente, conservando su orden de dependencia.
    \item \textit{Corrección de fallas o inconsistencias} se concentra en S23, la semana en que se congela la implementación, para que las campañas experimentales posteriores se ejecuten sobre una versión estable del sistema.
    \item \textit{Recolección de información} se ejecuta de forma simultánea a \textit{Ejecución de pruebas de evaluación}, dado que el registro de las métricas se realiza de manera automática durante cada corrida.
    \item \textit{Preparación de la sustentación final} se adelanta a S26--S28, en correspondencia con la fecha prevista de sustentación en octubre.
    \item \textit{Elaboración del documento final} se ubica a partir de S29, con posterioridad a la sustentación, conforme a la norma del espacio académico.
    \item \textit{Registro de avances del proyecto} conserva su ventana completa, sin modificación alguna.
\end{itemize}

Estos ajustes se compensan entre sí: el traslado de la elaboración del documento final a la etapa posterior a la sustentación libera las semanas que permiten absorber el desplazamiento de las actividades técnicas.

\subsection{Organización del trabajo}

Entre S18 y S23 el trabajo se distribuye en dos frentes simultáneos, que convergen a partir de S24 cuando las campañas experimentales y el análisis de resultados operan sobre un sistema único.

\begin{table}[H]
    \centering
    \small
    \caption{Frentes de trabajo.}
    \label{tab:frentes}
    \begin{tabular}{@{}cp{9.0cm}l@{}}
        \toprule
        \textbf{Frente} & \textbf{Contenido} & \textbf{Semanas} \\ \midrule
        A & Coordinación en simulación: réplica del entorno, dos agentes, nodo de coordinación, protocolo de relevo, instrumentación e interfaz de usuario & S17--S23 \\ \addlinespace
        B & Plataforma física: diagnóstico de hardware, paridad entre vehículos, mapeo del entorno real, puesta en marcha y despliegue & S18--S23 \\ \addlinespace
        T & Transversal: requisitos, protocolo experimental y registro de avances & S17--S32 \\ \bottomrule
    \end{tabular}
\end{table}

\section{Cronograma detallado}

Cada semana cierra con un \textbf{criterio verificable}: una semana no se da por cumplida por actividad realizada, sino por artefacto producido.

\subsection{S17 (3--9 ago) --- Consolidación del stack de navegación y contrato de interfaces}

\begin{table}[H]
    \centering
    \small
    \begin{tabular}{@{}cp{12.8cm}@{}}
        \toprule
        \textbf{Frente} & \textbf{Actividad} \\ \midrule
        A & Consolidar y publicar en el repositorio el stack de navegación autónoma \\
        B & Instalar el entorno de cómputo en la Raspberry Pi 4 y en la tarjeta original del vehículo \\
        B & Verificar el acceso por red a ambas tarjetas y la locomoción del vehículo \\
        A & Definir el contrato de interfaces ROS 2: espacios de nombres, marcos de referencia, acciones y tipos de mensaje \\
        T & Actualizar el tablero de trazabilidad y emitir los entregables correspondientes \\ \bottomrule
    \end{tabular}
\end{table}

\textbf{Criterio de cierre:} contrato de interfaces documentado; vehículo físico con locomoción verificada; entregables emitidos.

\subsection{S18 (10--16 ago) --- Réplica del entorno y diagnóstico de hardware}

\begin{table}[H]
    \centering
    \small
    \begin{tabular}{@{}cp{12.8cm}@{}}
        \toprule
        \textbf{Frente} & \textbf{Actividad} \\ \midrule
        A & Cerrar el contrato de interfaces: dos agentes con espacios de nombres independientes coexistiendo en el simulador \\ \addlinespace
        A & Extraer la planta \texttt{primer\_piso\_v2} como modelo reutilizable y replicarla en dos niveles con separación vertical y zona de transición señalizada en cada planta \\ \addlinespace
        A & Verificar navegación autónoma sobre el nivel superior \\ \addlinespace
        B & Diagnóstico acotado de hardware sobre cuatro puntos: lecturas utilizables del LiDAR real, comando del vehículo desde ROS 2 sin la interfaz web, latencia de comunicación entre los dos vehículos, e implicaciones de la distribución de ROS 2 empleada \\ \bottomrule
    \end{tabular}
\end{table}

\textbf{Criterio de cierre:} dos agentes coexisten con marcos y tópicos aislados; el mundo de dos niveles carga en el simulador y el vehículo navega sobre el nivel superior; informe del diagnóstico con los cuatro resultados.

\textbf{Nota (2026-08-16).} La redacción original de esta semana pedía además \textit{generar el mapa del nivel superior}. Esa actividad queda \textbf{suprimida}: las dos plantas son la misma geometría en las mismas coordenadas horizontales, de modo que un solo mapa sirve para ambos niveles, y un segundo mapa de la misma planta solo podría divergir del primero. No se suprime por argumento sino por medida: los cortes horizontales a 0,30 m y a 3,30 m de \texttt{primer\_piso\_dos\_niveles.world} producen mapas idénticos píxel a píxel, y ambos superan \texttt{verificar\_mapa.py} con 99,7 \% de cobertura y cero obstáculos sin pared real detrás. La navegación sobre el nivel superior sí se verificó, y esa parte de la actividad se conserva: 4,79 m de recorrido medidos contra \texttt{/odom} sobre un objetivo de 5,0 m, con la altura constante a 1,9 micrómetros.

\subsection{S19 (17--23 ago) --- Navegación por agente, requisitos y protocolo}

\begin{table}[H]
    \centering
    \small
    \begin{tabular}{@{}cp{12.8cm}@{}}
        \toprule
        \textbf{Frente} & \textbf{Actividad} \\ \midrule
        A & Configurar una instancia de navegación independiente por agente, una por nivel \\
        B & Llevar el segundo vehículo al mismo estado del primero \\
        B & Resolver los puntos que haya arrojado el diagnóstico de hardware \\
        T & Construir la matriz de requisitos que relaciona requisito funcional, objetivo específico y prueba de verificación \\
        T & Definir el protocolo experimental: qué se mide, cómo, cuántas repeticiones y con qué criterios de éxito, con la continuidad del servicio entre niveles como variable de respuesta principal \\ \bottomrule
    \end{tabular}
\end{table}

\textbf{Criterio de cierre:} cada agente ejecuta órdenes de navegación dentro de su nivel sin interferencia entre marcos de referencia ni entre tópicos. Matriz de requisitos publicada y protocolo experimental escrito antes de instrumentar. La matriz y el protocolo se elaboran juntos: un requisito sin prueba asociada no es un requisito.

\subsection{S20 (24--30 ago) --- Nodo de coordinación}

\begin{table}[H]
    \centering
    \small
    \begin{tabular}{@{}cp{12.8cm}@{}}
        \toprule
        \textbf{Frente} & \textbf{Actividad} \\ \midrule
        A & Servidor de coordinación: registro de agentes, recepción de la solicitud del usuario y asignación por nivel \\
        A & Instrumentar el tiempo de respuesta y el tiempo de asignación \\
        B & Mapear con el vehículo físico las dos plantas del pasillo y contrastar el mapa resultante contra el modelo \texttt{primer\_piso\_v2} \\ \bottomrule
    \end{tabular}
\end{table}

\textbf{Criterio de cierre:} una solicitud de origen y destino produce la asignación del agente correcto y deja registrados ambos tiempos. Mapas de las dos plantas guardados y verificados; la discrepancia entre el mapa levantado y el modelo simulado queda medida en metros.

\subsection{S21 (31 ago--6 sep) --- Protocolo de relevo}

\begin{table}[H]
    \centering
    \small
    \begin{tabular}{@{}cp{12.8cm}@{}}
        \toprule
        \textbf{Frente} & \textbf{Actividad} \\ \midrule
        A & Máquina de estados del relevo: guiado hasta la zona de transición, entrega, notificación al segundo agente, espera y reanudación \\
        A & Instrumentar la tasa de éxito y la continuidad del servicio \\
        B & Navegación autónoma punto a punto sobre el vehículo físico en el laboratorio \\ \bottomrule
    \end{tabular}
\end{table}

\textbf{Criterio de cierre:} en simulación, una solicitud que cruza entre niveles se completa de extremo a extremo con relevo, y el registro arroja las cuatro métricas de evaluación.

Esta semana concentra además el \textbf{punto de decisión} sobre el segundo vehículo. Si los dos vehículos quedan operativos, el relevo físico se demuestra con un vehículo por planta; en caso contrario, se demuestra con un vehículo real y uno simulado, y así se declara. En ninguno de los dos casos se modifican las fechas del cronograma ni se compromete la evidencia del relevo, que está cubierta por las pruebas en simulación.

\textbf{Criterio de habilitación de las pruebas en el pasillo:} se habilitan cuando un vehículo navegue punto a punto de forma autónoma en el laboratorio con localización estable. El relevo completo entre dos vehículos no es criterio de habilitación: es el núcleo del cumplimiento del proyecto, y condicionar el avance a su logro previo bloquearía la secuencia.

\subsection{S22 (7--13 sep) --- Interfaz de usuario e integración}

\begin{table}[H]
    \centering
    \small
    \begin{tabular}{@{}cp{12.8cm}@{}}
        \toprule
        \textbf{Frente} & \textbf{Actividad} \\ \midrule
        A & Interfaz de interacción humano--robot con selección de origen y destino \\
        A & Integrar los módulos del sistema de extremo a extremo en simulación \\
        B & Desplegar el sistema sobre los vehículos físicos \\ \bottomrule
    \end{tabular}
\end{table}

\textbf{Criterio de cierre:} desde un dispositivo móvil se selecciona origen y destino, y el sistema completa el guiado con relevo sin intervención manual.

\subsection{S23 (14--20 sep) --- Verificación, corrección y cierre de la implementación}

\begin{table}[H]
    \centering
    \small
    \begin{tabular}{@{}cp{12.8cm}@{}}
        \toprule
        \textbf{Frente} & \textbf{Actividad} \\ \midrule
        A+B & Verificación del funcionamiento general y corrección de fallas o inconsistencias \\
        B & Ensayo del protocolo en el pasillo con los vehículos reales \\
        T & Congelar la implementación: a partir de este punto no se incorpora funcionalidad nueva \\ \bottomrule
    \end{tabular}
\end{table}

\textbf{Criterio de cierre:} el sistema opera de extremo a extremo en simulación y se ejecuta al menos una corrida física completa. Repositorio etiquetado.

\subsection{S24 (21--27 sep) --- Pruebas de evaluación en simulación}

\begin{table}[H]
    \centering
    \small
    \begin{tabular}{@{}p{13.8cm}@{}}
        \toprule
        \textbf{Actividad} \\ \midrule
        Ejecutar 30 repeticiones del protocolo completo con registro automático de las cuatro métricas \\
        Consolidar y versionar el conjunto de datos resultante \\ \bottomrule
    \end{tabular}
\end{table}

\textbf{Criterio de cierre:} 30 corridas registradas y conjunto de datos versionado.

\subsection{S25 (28 sep--4 oct) --- Pruebas de evaluación en el entorno físico}

\begin{table}[H]
    \centering
    \small
    \begin{tabular}{@{}p{13.8cm}@{}}
        \toprule
        \textbf{Actividad} \\ \midrule
        Ejecutar entre 5 y 10 repeticiones en el pasillo con los vehículos reales, un vehículo por planta, en franjas de baja circulación \\
        Grabar y editar el video de la demostración \\ \bottomrule
    \end{tabular}
\end{table}

\textbf{Criterio de cierre:} corridas físicas registradas con la misma instrumentación empleada en simulación y sobre la misma geometría que se simuló; video de la demostración terminado.

\subsection{S26 (5--11 oct) --- Análisis de resultados}

\begin{table}[H]
    \centering
    \small
    \begin{tabular}{@{}p{13.8cm}@{}}
        \toprule
        \textbf{Actividad} \\ \midrule
        Análisis comparativo entre los resultados de simulación y los del entorno físico \\
        Verificación del sistema contra la matriz de requisitos \\
        Iniciar la preparación de la sustentación \\ \bottomrule
    \end{tabular}
\end{table}

\textbf{Criterio de cierre:} gráficas y tablas de las cuatro métricas; tabla de cumplimiento requisito por requisito.

\subsection{S27 (12--18 oct) --- Conclusiones y preparación de la sustentación}

\begin{table}[H]
    \centering
    \small
    \begin{tabular}{@{}p{13.8cm}@{}}
        \toprule
        \textbf{Actividad} \\ \midrule
        Elaborar las conclusiones sobre la viabilidad técnica de la solución propuesta \\
        Armar la presentación y ensayarla de forma cronometrada \\ \bottomrule
    \end{tabular}
\end{table}

\textbf{Criterio de cierre:} presentación completa y ensayada de principio a fin.

\subsection{S28 (19--25 oct) --- Sustentación}

\begin{table}[H]
    \centering
    \small
    \begin{tabular}{@{}p{13.8cm}@{}}
        \toprule
        \textbf{Actividad} \\ \midrule
        Ensayos finales y plan de contingencia de la demostración en vivo \\
        Sustentación \\ \bottomrule
    \end{tabular}
\end{table}

\textbf{Criterio de cierre:} sustentación realizada; observaciones del jurado registradas por escrito.

\subsection{S29 (26 oct--1 nov) --- Sustentación e inicio del documento final}

\begin{table}[H]
    \centering
    \small
    \begin{tabular}{@{}p{13.8cm}@{}}
        \toprule
        \textbf{Actividad} \\ \midrule
        Sustentación, en caso de que la fecha corresponda a esta semana \\
        Definir la estructura del documento final e iniciar su redacción \\ \bottomrule
    \end{tabular}
\end{table}

\subsection{S30--S31 (2--15 nov) --- Elaboración del documento final}

\begin{table}[H]
    \centering
    \small
    \begin{tabular}{@{}p{13.8cm}@{}}
        \toprule
        \textbf{Actividad} \\ \midrule
        Redacción completa del documento, incorporando las observaciones del jurado \\
        Revisión con los directores del proyecto \\ \bottomrule
    \end{tabular}
\end{table}

\textbf{Criterio de cierre (S31):} borrador completo entregado a los directores.

\subsection{S32 (16--22 nov) --- Entrega}

\begin{table}[H]
    \centering
    \small
    \begin{tabular}{@{}p{13.8cm}@{}}
        \toprule
        \textbf{Actividad} \\ \midrule
        Incorporar las correcciones de los directores \\
        Entrega final \\ \bottomrule
    \end{tabular}
\end{table}

\section{Hitos}

Los hitos señalan los puntos en que el proyecto adquiere una capacidad que habilita el trabajo posterior. Su función es de control: si un hito se retrasa, se identifica de inmediato qué actividades quedan condicionadas.

\begin{table}[H]
    \centering
    \small
    \caption{Hitos del periodo S17--S32.}
    \label{tab:hitos}
    \begin{tabular}{@{}cp{6.8cm}cp{4.2cm}@{}}
        \toprule
        \textbf{\#} & \textbf{Hito} & \textbf{Semana} & \textbf{Habilita} \\ \midrule
        H1 & Contrato de interfaces ROS 2 definido & S17 & Desarrollo multi-robot \\
        H2 & Riesgos de hardware caracterizados & S18 & Puesta en marcha del vehículo \\
        H3 & Dos agentes navegando en niveles separados & S19 & Nodo de coordinación \\
        H4 & Asignación dinámica de tareas funcionando & S20 & Protocolo de relevo \\
        H5 & Relevo completo con métricas en simulación & S21 & Pruebas de evaluación \\
        H6 & Sistema integrado de extremo a extremo & S22 & Verificación \\
        H7 & Implementación congelada & S23 & Pruebas de evaluación \\
        H8 & Conjunto de datos completo & S25 & Análisis de resultados \\
        H9 & Sustentación & S28--S29 & Documento final \\ \bottomrule
    \end{tabular}
\end{table}

\section{Seguimiento}

El avance se registra semanalmente en el tablero de trazabilidad del proyecto, donde se consignan el porcentaje de avance de cada objetivo específico, el estado de los riesgos abiertos y el criterio de cierre alcanzado o no alcanzado. Este procedimiento materializa la actividad \textit{Registro de avances del proyecto}, contemplada en el anteproyecto para la totalidad de las semanas del cronograma.

\begin{thebibliography}{99}

\bibitem{anteproyecto}
S. Hernández Ávila y J. A. Mejía León,
``Sistema colaborativo de robots móviles para asistencia de orientación en entornos interiores con múltiples pisos --- Anteproyecto,''
Universidad Santo Tomás, Facultad de Ingeniería Electrónica, Bogotá, 2026.

\bibitem{ros2_humble_docs}
Open Robotics,
``ROS 2 Humble Hawksbill Documentation,''
[Online]. Available: \url{https://docs.ros.org/en/humble/}

\bibitem{nav2_docs}
Open Robotics,
``Nav2 Documentation,''
[Online]. Available: \url{https://docs.nav2.org/}

\end{thebibliography}

\end{document}
